\sect{श्रावण-कृष्ण-कामिका-एकादशी-माहात्म्यम्}
\label{sec:padma-shravana-krishna-kamika}

\ganapatyadiDhyanam
\padmaPuranam

\dnsub{अथ षट्पञ्चाशत्तमोऽध्यायः॥५६}

\uvacha{युधिष्ठिर उवाच}

\twolineshloka
{श्रावणस्यासिते पक्षे किं नामैकादशी भवेत्}
{तन्नः कथय गोविन्द वासुदेव नमोऽस्तु ते}% ॥१॥

\uvacha{श्रीकृष्ण उवाच}

\twolineshloka
{शृणु राजन् प्रवक्ष्यामि आख्यानं पापनाशनम्}
{यत् प्रोक्तं ब्रह्मणा पूर्वं पृच्छते नारदाय वै}% ॥२॥

\uvacha{नारद उवाच}

\twolineshloka
{भगवन् श्रोतुमिच्छामि त्वत्तोऽहं कमलासन}
{श्रावणस्यासिते पक्षे किं नामैकादशी भवेत्}% ॥३॥

\twolineshloka
{को देवः को विधिस्तस्याः किं पुण्यं कथय प्रभो}
{इति तस्य वचः श्रुत्वा ब्रह्मा वचनमब्रवीत्}% ॥४॥

\uvacha{ब्रह्मोवाच}

\twolineshloka
{शृणु नारद ते वच्मि लोकानां हितकाम्यया}
{श्रावणैकादशी कृष्णा कामिका नाम नामतः}% ॥५॥

\twolineshloka
{अस्याः श्रवणमात्रेण वाजपेयफलं लभेत्}
{अस्यां यजति देवेशं शङ्खचक्रगदाधरम्}% ॥६॥

\twolineshloka
{श्रीधराख्यं हरिं विष्णुं माधवं मधुसूदनम्}
{पूजयेद् ध्यायते यो वै तस्य पुण्यफलं शृणु}% ॥७॥

\twolineshloka
{न गङ्गायां न काश्यां च नैमिषे न च पुष्करे}
{तत्फलं समवाप्नोति यत्फलं कृष्णपूजनात्}% ॥८॥

\twolineshloka
{गोदावर्यां गुरौ सिंहे व्यतीपाते च दण्डके}
{यत्फलं समवाप्नोति तत्फलं कृष्णपूजनात्}% ॥९॥

\twolineshloka
{ससागरवनोपेतां यो ददाति वसुन्धराम्}
{कामिकाव्रतकारी च ह्युभौ समफलौ स्मृतौ}% ॥१०॥

\twolineshloka
{प्रसूयमानां यो धेनुं दद्यात् सोपस्करां नरः}
{तत्फलं समवाप्नोति कामिकाव्रतकारकः}% ॥११॥

\twolineshloka
{श्रावणे श्रीधरं देवं पूजयेद्यो नरोत्तमः}
{तेनैव पूजिता देवा गन्धर्वोरगपन्नगाः}% ॥१२॥

\twolineshloka
{तस्मात् सर्वप्रयत्नेन कामिकादिवसे हरिः}
{पूजनीयो यथाशक्ति मानुषैः पापभीरुभिः}% ॥१३॥

\twolineshloka
{ये संसारार्णवे मग्नाः पापपङ्कसमाकुले}
{तेषामुद्धरणार्थाय कामिकाव्रतमुत्तमम्}% ॥१४॥

\twolineshloka
{नातः परतरा काचित् पवित्रा पापहारिणी}
{एवं नारद जानीहि स्वयमाह परो हरिः}% ॥१५॥

\twolineshloka
{अध्यात्मविद्या निरतैर्यत्फलं प्राप्यते नरैः}
{ततो बहुतरं विद्धि कामिकाव्रतसेविनाम्}% ॥१६॥

\twolineshloka
{रात्रौ जागरणं कृत्वा कामिकाव्रतकृन्नरः}
{न पश्यति यमं रौद्रं नैव गच्छति दुर्गतिम्}% ॥१७॥

\twolineshloka
{न पश्यति कुयोनिं च कामिकाव्रतसेवनात्}
{कामिकाया व्रते चीर्णे कैवल्यं योगिनो गताः}% ॥१८॥

\twolineshloka
{तस्मात् सर्वप्रयत्नेन कर्तव्या नियतात्मभिः}
{तुलसीप्रभवैः पत्रैः यो नरः पूजयेद्धरिम्}% ॥१९॥

\twolineshloka
{न लिप्यते स पापेन पद्मपत्रमिवाम्भसा}
{सुवर्णभारमेकं तु रजतं च चतुर्गुणम्}% ॥२०॥

\twolineshloka
{तत्फलं समवाप्नोति तुलसीदलपूजनात्}
{रक्त-मौक्तिक-वैदूर्य-प्रवालादिभिरर्चितः}% ॥२१॥

\twolineshloka
{न तुष्यति तथा विष्णुस्तुलसीदलतो यथा}
{तुलसीमञ्जरीभिश्च पूजितो येन केशवः}% ॥२२॥

\fourlineindentedshloka
{या दृष्टा निखिलाघसङ्घशमनी स्पृष्टा वपुः पावनी}
{रोगाणामभिवन्दितानि रसिनी सिक्तान्तकत्रासिनी}
{प्रत्यासत्तिविधायिनी भगवतः कृष्णस्य संरोपिता}
{न्यस्ता तच्चरणे विमुक्तिफलदा तस्यै तुलस्यै नमः}% ॥२३॥

\twolineshloka
{दीपं ददाति यो मर्त्यो दिवारात्रं हरेर्दिने}
{तस्य पुण्यस्य सङ्ख्यातुं चित्रगुप्तो न वेत्त्यलम्}% ॥२४॥

\twolineshloka
{कृष्णाग्रे दीपको यस्य ज्वलत्येकादशीदिने}
{पितरस्तस्य तृप्यन्ति अमृतेन दिवि स्थिताः}% ॥२५॥

\twolineshloka
{घृतेन दीपं प्रज्वाल्य तिलतैलेन वा पुनः}
{प्रयाति सूर्यलोकं च दीपकोटिशतार्चितः}% ॥२६॥

\twolineshloka
{अयं तवाग्रे कथितः कामिकामहिमा मया}
{अतो नरैः प्रकर्तव्या सर्वपातकहारिणी}% ॥२७॥

\twolineshloka
{ब्रह्महत्यापहरणी भ्रूणहत्याविनाशिनी}
{वैष्णवस्थानदात्री च महापुण्यफलप्रदा}% ॥२८॥

\twolineshloka
{श्रुत्वा माहात्म्यमेतस्या नरः श्रद्धासमन्वितः}
{विष्णुलोकमवाप्नोति सर्वपापैः प्रमुच्यते}% ॥२९॥

॥इति श्रीपाद्मे महापुराणे पञ्चपञ्चाशत्साहस्र्यां संहितायामुत्तरखण्डे उमापतिनारदसंवादान्तर्गतकृष्णयुधिष्ठिरसंवादे श्रावण-कृष्ण-कामिका-एकादशी-माहात्म्यं नाम षट्पञ्चाशत्तमोऽध्यायः॥५६॥

\genMangalaShloka

\hyperref[sec:ekadashi_mahatmyam_padma_puranam]{\closesub}
\clearpage

